\documentclass[11pt]{article}
\usepackage{fontspec}
\usepackage[french]{babel}
\usepackage{amsmath}
\usepackage{amssymb}
\usepackage{amsthm}
\usepackage{mathrsfs}
\usepackage{tikz-cd}
\usepackage{enumitem}

\newcommand{\printedpage}[1]{\marginpar{\footnotesize p.~#1}}
\newcommand{\msp}[1]{\marginpar{\footnotesize [#1]}}
\newcommand{\circled}[1]{\tikz[baseline=(c.base)]\node[draw,circle,inner sep=1pt](c){#1};}
\newcommand{\Fib}{\underline{\mathrm{Fib}}\,}
\newcommand{\Fibsc}{\underline{\mathrm{Fibsc}}\,}
\newcommand{\Cofib}{\underline{\mathrm{Cofib}}\,}
\newcommand{\Cofibsc}{\underline{\mathrm{Cofibsc}}\,}
\newcommand{\Fibhot}{\underline{\mathrm{Fibhot}}\,}
\newcommand{\Parf}{\underline{\mathrm{Parf}}\,}
\newcommand{\Liss}{\underline{\mathrm{Liss}}}
\newcommand{\Prop}{\underline{\mathrm{Prop}}}
\newcommand{\Ho}{\mathrm{Ho}}
\newcommand{\Hom}{\mathrm{Hom}}
\newcommand{\HOT}{\mathrm{HOT}}
\newcommand{\hot}{\mathrm{hot}}
\newcommand{\Fl}{\mathrm{Fl}}
\newcommand{\Cat}{\mathrm{Cat}}
\newcommand{\M}{\mathcal{M}}

\begin{document}
\setcounter{footnote}{59}

\printedpage{69}
% head: A. Grothendieck : Dérivateurs, Ch. VI, Édition des Universités de Montpellier II et Paris VII
qui est donné sur les objets par
\[\gamma(X) = C(X),\]
et qui sur les flèches se définit de façon évidente dans le cas A) (fonctorialité de $C(X)$ en $X$), et dans le cas de B) se construit ad hoc, comme la seule loi fonctorielle pour laquelle les homomorphismes
\[i_X : X \longrightarrow C(X)\]
\msp{page 110}
(interprétés comme flèches dans $(\M/S)W_S^{-1}$) soient eux-mêmes fonctoriels en $X$.

On désigne par
\[\Ho^{\mathrm{lc}}_W(S) \qquad\qquad \text{(lc comme ‘localement constant’)}\]
la catégorie des fractions $(\M/S)W_S^{-1}$, équivalente à $(\Fib S)W_S^{-1}$.

\circled{II} Donnons nous un localiseur $W_S^0$ dans $\M/S$
\[W_S^0 \subseteq W_S \subseteq \Fl(\M/S),\]
avec l'hypothèse formelle similaire à celle de $W$ (cf. 1$^o$). On suppose que

\begin{itemize}[leftmargin=4em]
\item[3$^o$)] $W_S^0 \cap \Fl(\Fib S) = W_S|\Fib S$.
\end{itemize}

(Cf. pages 94-96 pour une façon générale d'obtenir des tels localiseurs. Il faut bien voir que pour définir des ‘types d'homotopie’ relatifs sur $S$ le localiseur $W_S$ est beaucoup trop grossier, moralement il ne donne qu'une catégorie des types d'homotopie relatifs localement constants, et c'est justement ça qui est montré par l'equivalence de \circled{I} , où $\Fib S$ est visualisé comme formé des ‘fibrations localement triviales’ sur $S$.) On a alors un deuxième localisé % [sic] l'equivalence -> l'équivalence, accent manquant
\[\Ho_{W_S^0}(S) \overset{\text{déf}}{=} (\M/S)(W_S^0)^{-1},\]
et deux foncteurs
\[
\begin{tikzcd}[column sep=large]
\Ho_{W_S^0}(S) \arrow[r, shift left, "\rho"] & \Ho^{\mathrm{lc}}_W(S), \arrow[l, shift left, "i"]
\end{tikzcd}
\tag*{$(*)$}
\]
où $\rho$ est le foncteur de localisation (strict) évident, dû à l'inclusion $W_S^0 \subseteq W_S$, et où $i$ est
\msp{page 111}
essentiellement le foncteur déduit de l'inclusion
\[(\Fib S, \underbrace{W_S|\Fib S}_{=W_S^0|\Fib S}) \longrightarrow (\M/S, W_S^0)\]
par localisation \textit{au moyen de $W_S^0$}, donc le foncteur
\[(\Fib S)(W_S^0)^{-1} \longrightarrow (\M/S)(W_S^0)^{-1}, \tag*{$(**)$}\]
% foot: Version 9 novembre 2001
\printedpage{70}
% head: A. Grothendieck : Dérivateurs, Ch. VI, Édition des Universités de Montpellier II et Paris VII
compte tenu de l'équivalence $(*)$ de \circled{I} , i.e. $i$ est le composé du quasi-inverse de $(*)$, avec le foncteur $(**)$ précédent. Explicitement, il est défini par
\[i(X) = \text{classe de C(X) dans } (\M/S)(W_S^0)^{-1},\]
en tenant compte du fait que $X \longmapsto C(X)$ transforme $W_S$ en $W_S^o$, grâce à la condition 3$^o$) sur $W_S^0$. % [sic] W_S^o -> W_S^0, lettre o au lieu du chiffre 0 utilisé partout ailleurs

Ceci posé, \textit{sous l'hypothèse} B (‘Quillen’), on a : dans le couple de foncteurs $(\rho, i)$ de $(*)$, page 110, \textit{$i$ est adjoint à droite de $\rho$,} i.e. on a pour $\xi \in \mathrm{Ob}\,\Ho_{W_S^0}(S)$, $\eta \in \mathrm{Ob}\,\Ho^{\mathrm{lc}}_W(S)$
\[\Hom_{W_S^0}(\xi, i(\eta)) \simeq \Hom_{W_S}(\rho(\xi), \eta),\]
ce qui s'explicite, pour $X$, $Y$ dans $\M/S$, par
\[\boxed{\Hom_{W_S^0}(X, C(Y)) \simeq \Hom_{W_S}(X, Y)}\]
(\footnote{\textbf{NB} Le premier membre est indépendant du choix de $W_S^0$! (Pourvu seulement que l'hypothèse $W_S^0 \cap \Fl\,\Fib S = W_S \cap \Fl\,\Fib S$, en plus de $W_S^0 \subseteq W_S$, soit satisfaite.)}). Il s'ensuit que \textit{le foncteur $i$ est pleinement fidèle.}

La même formule d'adjonction, et la
\msp{page 112}
pleine fidélité de $i$, sont valables dans le cas A, quand on renforce l'hypothèse A (relative à un $S$ particulier) en une hypothèse plus restrictive et plus compréhensive, que voici.

\begin{itemize}[leftmargin=3em]
\item[A$'$)] Le foncteur $C : \M/S \longrightarrow \M/S$ et le morphisme fonctoriel $i_X : X \longrightarrow C(X)$ sont construits ‘à la \textsc{Cartan-Serre}’ à partir d'un foncteur défini sur $\M$ entier
\[\M \xrightarrow{\mathrm{Ch}} \M,\]
du type ‘chemins’, i.e. muni en plus des données suivantes
\begin{itemize}[leftmargin=3em]
\item[a)] \[X \xrightarrow{\alpha_X} \mathrm{Ch}(X) \xrightarrow{p_X = (\sigma_X, \beta_X)} X \times X\]
(structure de source et but des chemins, et de chemins triviaux), et
\item[b)] Une composition des chemins, sans exiger ici l'associativité de la composition mais seulement que les chemins triviaux soient unités bilatères pour la composition. De plus, on suppose que pout tout $X$ % [sic] pout -> pour, faute de frappe
\[\begin{cases} \alpha_X \in W,\ p_X \in \mathrm{Fib} \\ (X \longrightarrow e) \in \mathrm{Fib}. \end{cases}\]
\end{itemize}
\end{itemize}
% foot: Version 9 novembre 2001

\printedpage{71}
% head: A. Grothendieck : Dérivateurs, Ch. VI, Édition des Universités de Montpellier II et Paris VII
\msp{page 113}
\textbf{Corollaire.} Soient $X$ dans $\M/S$, $Y$ dans $\Fib S$, alors
\[\Hom_{W_S^0}(X, Y) \xrightarrow{\sim} \Hom_{W_S}(X, Y).\]
\textsc{Démonstration}. Par la formule d'adjonction (p. 111), le deuxième membre est $\Hom_{W_S^0}(X, CY)$. Mais comme $Y$ et $C(Y)$ sont dans $\Fib S$, où $W_S$ et $W_S^0$ coïncident, $Y \xrightarrow{i_Y} C(Y)$ est non seulement dans $W_S$, mais dans $W_S^0$, donc $\Hom_{W_S^0}(X, Y) \xrightarrow{\sim} \Hom_{W_S^0}(X, CY)$, d'où le résultat.

Revenons à Cat, Cat$/S$. La digression précédente sur $(\M, W, \mathrm{Fib})$ a permis de mieux situer ce qui concerne le foncteur $C$, et les relations entre les catégories localisées de
\[\Cat/S \ \ (\text{pour } W_S,\ W_S^{\mathrm{d}},\ W_S^{\mathrm{g}}) \qquad\quad \text{et} \qquad\quad \Parf S \ \ (\text{pour } W_S^{\mathrm{f}} = W_S = W_S^{\mathrm{d}} = W_S^{\mathrm{g}})\]
ainsi que pour les catégories intermédiaires
\[\Parf S \subseteq \mathcal{F} \subseteq \Cat/S\]
pleines de Cat$/S$ (sur le diagramme p. 76, il y en a cinq autres, savoir $\Fibhot S$, $\Liss\,S$, $\Liss_0 S$, $\Prop\,S$, $\Prop_0\,S$). Pour ce qui concerne les relations entre les localisés de ces sept catégories, qui se trouvent au centre du diagramme, les énoncés du théorème-scholie 1 (p. 70-81) ne donnent rien de plus que le contenu du théorème-scholie 2 (p. 108-113). Pour les établir, on a donc le choix entre les deux hypothèses A) ou B) dudit théorème. Soit utiliser un foncteur $C_S$ explicite, à l'aide d'un
\msp{page 114}
foncteur $X \longmapsto \underline{\mathrm{Ch}}(X)$ (‘chemins dans $X$’) sur Cat, donc les structures de source, but et composition pour les chemins, et les chemins triviaux (agissant comme unités bilatères pour la composition). La construction d'un tel foncteur $C$ n'a par l'air trivial et semble ne pas être connue - j'ai peiné pour y arriver! (Avec, en plus, l'associativité de la composition \ldots). Ou bien utiliser, à la \textsc{Quillen}, l'existence d'une classe de morphismes dits Q-fibrants, avec des bonnes propriétés de factorisation. Ce n'est pas non plus évident dans Cat, et c'est établi par \textsc{Thomason}, par une voie un peu compliquée. Peut-être serai-je amené à revenir dessus - mais seulement en cas de besoin! % [sic] n'a par l'air -> n'a pas l'air, faute de frappe

Par contre, les assertions concernant les catégories $\Fib S$, $\Fibsc S$, ou $\Cofib S$, $\Cofibsc S$, et leurs variantes à indices 0 (donc les huit catégories-ailes du diagramme), et leurs relations à Cat$/S$ d'une part, à $\Fibhot S$ de l'autre (pour les catégories à indice 0), ne peuvent être contenues dans un tapis général sur des catégories de modèles $(\M, W)$, me semble-t-il, vu la nature très spéciale des quatre catégories désignées ci-dessus, de leur rôle très particulier dans le formalisme catégorique. Or tout au moins les catégories ‘scindées’ $\Fibsc S$, $\Cofibsc S$, et par suite aussi leurs variantes à indice 0, jouent un rôle essentiel pour la théorie des HOT$(S)$, vu que c'est elles qui
\msp{page 115}
interviennent dualement quand on applique la recette générale de construction d'un dérivateur à partir d'une catégorie de modèles $(\M, W)$. Donc je ne vois pas de simplification
% foot: Version 9 novembre 2001
\printedpage{72}
% head: A. Grothendieck : Dérivateurs, Ch. VI, Édition des Universités de Montpellier II et Paris VII
substantielle à apporter au théorème-scholie 1. Tout au plus pourrait-on élaguer le diagramme de catégories envisagées, en laissant tomber $\Liss\,S$, $\Liss_0 S$, $\Prop\,S$, $\Prop_0 S$, qui n'y jouent pas un rôle vraiment utile. Pour ce qui concerne les résultats énoncés, ce qui les concerne est couvert par ce qui est dit pour les sous-catégories pleines $\mathcal{F}$ de Cat$/S$ qui contiennent (disons) $\Parf S$ (ou au choix celles qui contiennent $\Fib S$ ou $\Cofib S$).

La description des catégories du type HOT$(S)$, pour $S$ fixé, me paraît à présent parfaitement bien comprise. Il est temps de faire varier $S$, et de comprendre le caractère contravariant en $S$. On se donne donc
\[S' \xrightarrow{f} S,\]
et on se propose de décrire les foncteurs image inverse pour les trois types de catégories HOT$(S)$. Mais nous allons nous borner à présent aux $\HOT_W(S)$ et $\HOT^{\mathrm{lc}}_W(S)$, puisque tout ce qui concerne $\HOT^o_W(S) = \HOT_W(S^o)$ \texttt{[plutôt} $\HOT_W(S^o)^o$\texttt{, cf. p. 50]} s'en déduira par dualité $S \longmapsto S^o$. Donc il s'agit de décrire
\msp{page 116}
\[f^* : \HOT_W(S) \longrightarrow \HOT_W(S') \tag*{$(*)$}\]
et le foncteur induit
\[f^* : \HOT^{\mathrm{lc}}_W(S) \longrightarrow \HOT^{\mathrm{lc}}_W(S').\]
Quand on décrit les catégories HOT$(S)$, conformément à la recette générale, à l'aide de foncteurs
\[F : S^o \longrightarrow \Cat\]
(préfaisceaux en catégories sur $S$), ou, ce qui revient au même, de catégories scindées, la description de la flèche $(*)$ est évidente : elle revient à la loi
\[\left|\begin{array}{rcl}
F & \longmapsto & F \circ f^o \\[1ex]
\underline{\mathrm{Hom}}(S^o, \Cat) & \longrightarrow & \underline{\mathrm{Hom}}(S'^o, \Cat),
\end{array}\right.\]
en passant aux localisés pour la $W$-équivalence argument par argument. Dans le langage des catégories fibrées, il est bien connu que cette opération triviale sur les préfaisceaux en catégories correspond à l'opération \textit{image inverse} de catégories sur $S$ - laquelle opération passe aux scindages. Cela nous donne donc
\[f^*(\hot_{W,S}(X)) = \hot_{W,S'}(\underbrace{X \times_S S'}_{X'}),\]
pour toute $X$ fibrée scindée sur $S$. Or on sait que le scindage de $X$ ne joue aucun rôle, finalement, dans la description de la classe
\[\hot_{W,S}(X) \in \HOT_W(S)\]
% foot: Version 9 novembre 2001
\printedpage{73}
% head: A. Grothendieck : Dérivateurs, Ch. VI, Édition des Universités de Montpellier II et Paris VII
\msp{page 117}
associée à $X$. On peut oublier le scindage, et même (pour étudier les $\Hom_{\mathrm{Hot}_W(S)}(X, Y)$) travailler avec des $S$-foncteurs \textit{quelconque} entre catégories fibrées sur $S$ - pas la peine même de les supposer cartésiens, voir même compatibles avec des scindages. Cela me paraît une chose très importante pour le formalisme des $\HOT_W(S)$, et très spéciale au cas de la catégorie de modèles Cat. C'est un peu tacite dans le théorème-scholie - mais mes doutes initiaux concernant la catégorie $\underline{\mathrm{Fibsc}}$ elle-même montrent bien que je sentais qu'il y avait là quelque chose de pas évident, ni même (presque) \texttt{[?]} de vraisemblable. Et pourtant l'existence du foncteur $\Phi$ montre qu'il en est bien ainsi.

On peut donc interpréter $f^*$ en termes de $\Fib S$, $\Fib S'$, par le foncteur changement de base
\[f^* : \Fib S \longrightarrow \Fib S', \qquad\quad X \longmapsto X \times_S S'/S',\]
lequel foncteur préserve évidemment la $W$-équivalence par fibres $W_S^{\mathrm{f}}$, $W_{S'}^{\mathrm{f}}$, qu'est celle pertinente ici
\[f^*(W_S^{\mathrm{f}}) \subseteq W_{S'}^{\mathrm{f}},\]
de sorte que $f^*$ induit bien
\[(\Fib S)(W_S^{\mathrm{f}})^{-1} \longrightarrow (\Fib S')(W_S^{\mathrm{f}})^{-1}\]
\texttt{[plutôt} $(\Fib S')(W_{S'}^{\mathrm{f}})^{-1}$\texttt{]},
\msp{page 118}
i.e. le foncteur $(*)$
\[f^* : \HOT_W(S) \longrightarrow \HOT_W(S').\]
D'autre part, $\HOT_W(S)$ peut être décrit aussi comme localisé de Cat$/S$ tout entier, par rapport au localiseur $W_S^{\mathrm{d}}$, lequel sur $\Fib S$, et même sur $\Liss\,S$, coïncide avec $W_S^{\mathrm{f}}$. Mais on se bute ici à la difficulté que le foncteur image inverse
\[f^* : \Cat/S \longrightarrow \Cat/S', \qquad\quad X \longmapsto X \times_S S'/S',\]
n'est \textit{pas} compatible en général (si $S$ n'est isomorphe à la catégorie ponctuelle, pour $S'$ quelconque) avec les localiseurs $W_S^{\mathrm{d}}$, $W_{S'}^{\mathrm{d}}$. On le voit p.ex. en prenant $S' = \{s\}$, et en se rappelant que $W_S^{\mathrm{d}}$ n'implique \textit{pas}, sur Cat$/S$ tout entier, la $W$-équivalence par fibres.

Ces ennuis n'ont pas lieu sûr $\Liss\,S$, où $W_S^{\mathrm{d}}$ et $W_S^{\mathrm{f}}$ coïncident. On peut donc décrire $f^*$ sur les HOT aussi à l'aide de % [sic] sûr -> sur, accent fautif
\[f^* : \Liss\,S \longrightarrow \Liss\,S', \qquad\quad X \longmapsto X \times_S S'/S',\]
en passant aux localisés
\[f^* : (\Liss\,S)(W_S^{\mathrm{f}})^{-1} \longrightarrow (\Liss\,S')(W_{S'}^{\mathrm{f}})^{-1}.\]
En d'autres termes :

\textbf{Proposition 1.} \textit{Si $X$ dans} Cat$/S$ \textit{est lisse sur $S$, alors on a un isomorphisme canonique}
\[f^*(\hot_{W,S}(X)) \simeq \hot_{W,S'}(X'), \qquad\quad \text{\textit{où} } X' = X \times_S S'/S'.\]
% foot: Version 9 novembre 2001

\end{document}
